% section: 特殊関数

\section{特殊関数}

1つ面白い話をしよう.
まず,次の3つの漸化式を見てほしい.
\[
  a_{n+1}=a_n^2,
  \quad
  a_{n+1}=a_n^2-1,
  \quad
  a_{n+1}=a_n^2-2.
\]
1つ目はきちんとここまで本書を読んできた諸君らならば,用意に解くことができるだろう.
実際,これは対数型漸化式であるから,両辺の対数を取って等比数列に帰着させることができる.

2つ目はどうだろうか$a_n$に指数がついているから,対数を使って処理したいが$-1$の項が邪魔していてうまくいかない.
不動点を見つけようにも指数の項が邪魔をするし,うまくいきそうにない.
少々意地悪であったかもしれないが,これは初等的な解法を用いては,その一般項を得られない漸化式である.

3つ目も一見先の例のように一般項が得難い見た目をしている.
しかしながらこれはある特殊な条件の噛み合いによって偶然にも初等的な解法が存在する.

ここまで,等差型,等比型,階差型,階比型など,漸化式を一定の型に
帰着して解く方法を学んできた.しかし,右辺が少し複雑になると,
同じ道具が急に働かなくなる.

ここで扱うのは,一般の漸化式がいつでも解けるという話ではない.
一見すると解けそうにない式の中にの数学的対象と偶然一致するものがある.
その一致を見抜くことで,初等的な型に還元できない漸化式を解いていく.
本章はそういう話である.

\subsection{基本対称式型}
導入に続いて,先程の漸化式を考えていこう.
\[
  a_1=3, \quad a_{n+1}=a_n^2-2
\]
この漸化式を解くに当たって邪魔をしているのは$-2$の項である.
これをうまく消せるような操作を考えたい.
また,\,2乗があることからうまく因数分解できないか考えてみよう.

以上の$2$と打ち消し会えるように,\,$2$が生まれるような対称式を考える.

2乗の項は正だから,
\[
  (a+b)^2=a^2+2ab+b^2
\]
で言うところの$2ab$が$2$になってくれるように調整すればうまくいきそうだ.
そのようなふうに考えていくと,
\[
  \left(\alpha+\frac{1}{\alpha}\right)^2=\alpha^2+2+\frac{1}{\alpha^2}
\]
という式が思い当たる.
実際,ここに$-2$をすると
\[
  \alpha^2+\frac{1}{\alpha^2}
\]
が残るから,成り立ちそうなことが分かる.
この実験を元に解答に移る.

\begin{quote}
\textbf{【問題】}
\begin{quote}
  次の漸化式で定義される数列の一般項を求めよ.
  \[
    a_1=3, \quad a_{n+1}=a_n^2-2
  \]
\end{quote}
\textbf{【解答】}\\一般項が
  \[
    b_n=\alpha^{2^{n-1}}+\alpha^{-2^{n-1}}
  \]
  で表される数列$b_n$を考える.
  \begin{align*}
    b_{n+1}&=\alpha^{2^{n}}-\alpha^{-2^{n}}\\
    &=\left(\alpha^{2^{n-1}}\right)^2-\left(\alpha^{-2^{n-1}}\right)^2\\
    &=\left(\alpha^{2^{n-1}}+\alpha^{-2^{n-1}}\right)^2-2\\
    &=b_n^2-2
  \end{align*}
  これは与式の漸化式と一致する.また,
  \[
    b_1=a_1=\alpha+\alpha^{-1}
  \]
  となる$\alpha$を考えると,
  \begin{align*}
    &\alpha+\alpha^{-1}=3\\
    &\alpha^2-3\alpha+1=0\\
    &\therefore \quad \alpha=\frac{3\pm \sqrt{5}}{2}
  \end{align*}
  得られた$\alpha$の2解は逆数の関係にある.
  以上より,
  \[
    a_n=\left(\frac{3+ \sqrt{5}}{2}\right)^{2^{n-1}}+\left(\frac{3- \sqrt{5}}{2}\right)^{2^{n-1}}
  \]
\begin{quote}
  
\end{quote}
\end{quote}

ここで,定数を $-2$ から $-1$ に変えてみよう.
\[
  a_{n+1}=a_n^2-1.
\]
同じ置換を試すと,
\begin{align*}
  a_n^2-1
    &=\left(u+u^{-1}\right)^2-1\\
    &=u^2+u^{-2}+1,
\end{align*}
となる.しかし,次の項として欲しい形は $u^2+u^{-2}$ である.
余分な $1$ が残るため,同じ対称式の族の中に戻ってこない.

漸化式の $-2$ と $-1$ は,見た目には小さな違いである.
しかし,対称式を保つ恒等式にとっては,その違いが決定的になる.
「ほとんど同じ式だから,ほとんど同じ方法で解ける」とは限らない.
解法が成立するかどうかは,係数の大きさではなく,恒等式が完全に一致
するかどうかで決まる.

同じ考え方は,三次式にも現れる.
\[
  \left(\alpha+\frac{1}{\alpha}\right)^3
  -3\left(\alpha+\frac{1}{\alpha}\right)
  =\alpha^3+\frac{1}{\alpha^3}.
\]
このように,対称な組み合わせを選ぶと,べき乗の複雑な展開が
再び同じ種類の式へ戻る.ただし,恒等式が少しでもずれると,その閉じた
構造は失われる.これが基本対称式型の面白さである.

これを発想にもつ漸化式を例題としよう.

\noindent\textbf{【問題】}\\
$a_1=\frac{8}{3},\quad a_{n+1}=a_n^3+3a_n$

\noindent\textbf{【解答】}\\
$a_n=b_n-\frac{1}{b_n}$\\
と置くと.

$\displaystyle \alpha-\alpha^{-1}=\frac{8}{3}$
を解くと,
$\displaystyle \alpha=-\frac{1}{3},\,3$






\subsection{三角関数型}
またもや似た見た目の漸化式を考える.
\[
  a_1=\frac{\sqrt{3}}{2},\quad a_{n+1}=2a_n^2-1
\]

この漸化式も前節と同様に,今までの方法ではうまく解くことができない.

今回の鍵は三角関数である.
三角関数と置くことで,数列の漸化式が角度の漸化式へ変わることがある.
その仕組みを確認するために,まず加法定理について復習しよう.
\begin{align*}
  \sin(\alpha+\beta)
    &=\sin\alpha\cos\beta+\cos\alpha\sin\beta\\
  \cos(\alpha+\beta)
    &=\cos\alpha\cos\beta-\sin\alpha\sin\beta\\
  \tan(\alpha+\beta)
    &=\frac{\tan\alpha+\tan\beta}
            {1-\tan\alpha\tan\beta}
\end{align*}
ここで $\alpha=\beta=x$ とすれば,倍角公式
\begin{align*}
  \sin 2x&=2\sin x\cos x\\
  \cos 2x&=2\cos^2x-1=1-2\sin^2x\\
  \tan 2x&=\frac{2\tan x}{1-\tan^2x}
\end{align*}
を得る.また,倍角公式から
\begin{align*}
  \sin^2x&=\frac{1-\cos 2x}{2}\\
  \cos^2x&=\frac{1+\cos 2x}{2}\\
  \tan^2x&=\frac{1-\cos 2x}{1+\cos 2x}
\end{align*}
を導くことができる.また,加法定理から,\,$3x=2x+1$とすることで,
\begin{align*}
  \sin 3x&=3\sin x-4\sin^3x\\
  \cos 3x&=4\cos^3x-3\cos x
\end{align*}
が求められる.

これらを踏まえて次の漸化式を考える.
\[
  a_1=\frac{\sqrt{3}}{2},\quad a_{n+1}=2a_n^2-1.
\]
右辺は $\cos 2x=2\cos^2x-1$ と一致している.さらに,
初項は
\[
  a_1=\cos\frac{\pi}{6}
\]
と書ける.そこで $a_n=\cos\theta_n$ と置くと,
\[
  \theta_1=\frac{\pi}{6},
  \quad
  \theta_{n+1}=2\theta_n
\]
と考えられる.従って候補は
\[
  a_n=\cos\left(2^{n-1}\frac{\pi}{6}\right)
\]
である.

この候補を帰納法で確認する.\,$n=1$ では初項に一致する.ある $n$ で
\[
  a_n=\cos\left(2^{n-1}\frac{\pi}{6}\right)
\]
が成り立つとすると,
\begin{align*}
  a_{n+1}
    &=2a_n^2-1\\
    &=2\cos^2\left(2^{n-1}\frac{\pi}{6}\right)-1\\
    &=\cos\left(2^n\frac{\pi}{6}\right).
\end{align*}
これは $n+1$ の式である.よって,
\[
  a_n=\cos\left(2^{n-1}\frac{\pi}{6}\right)
\]
が示された.

この解法の本質は,二次式を無理に展開したり因数分解したりすることではない.
数列の値を角度の関数として表すことで,二次式の計算を角度の2倍という単純な操作に翻訳しているのである.

今度は
\[
  a_1=\tan\frac{\pi}{12}=2-\sqrt{3},
  \quad
  a_{n+1}=\frac{2a_n}{1-a_n^2}
\]
を考える.右辺は$\tan$の倍角公式そのものである.したがって,
\[
  a_n=\tan\left(2^{n-1}\frac{\pi}{12}\right)
\]
と予想できる.

実際,\,$n=1$ で初項に一致し,帰納法の段階では
\begin{align*}
  a_{n+1}
    &=\frac{2a_n}{1-a_n^2}\\
    &=\cfrac{2\tan\left(2^{n-1}\cfrac{\pi}{12}\right)}
            {1-\tan^2\left(2^{n-1}\cfrac{\pi}{12}\right)}\\
    &=\tan\left(2^n\frac{\pi}{12}\right)
\end{align*}
となる.ただし,分母が $0$ になると$\tan$は定義できない.
三角関数型では,置換を思いつくだけでなく,途中で定義できなくなる項がないかも確認しなければならない.

余弦の例と$\tan$の例は,同じ方針を共有している.すなわち,
\[
  a_n=\varphi(\theta_n),
  \quad
  \varphi(m\theta)=F_m(\varphi(\theta))
\]
となる関数 $\varphi$ と整数倍公式を探し,角度の漸化式を解いている.
いわば媒介変数表示や置換のようなことをしている.

また,ここで現れた多項式
\[
  2x^2-1,
  \quad
  4x^3-3x
\]
は偶然の産物ではない.チェビシェフ多項式 $T_m(x)$ を
\[
  T_m(\cos\theta)=\cos(m\theta)
\]
で定義すれば,
\[
  T_0(x)=1,
  \quad
  T_1(x)=x,
  \quad
  T_{m+1}(x)=2xT_m(x)-T_{m-1}(x)
\]
となる.例えば $T_2(x)=2x^2-1$ である.
したがって,余弦の例は
\[
  a_{n+1}=T_2(a_n),
  \quad
  a_n=T_{2^{n-1}}(a_1)
\]
と書き直せる.三角関数の公式から始めて多項式の族へ進む道が,
付録のチェビシェフ多項式へつながっている.


\subsection{ニュートン法型}

少々離れたところから話を始める.
ある関数$f(x)$について,\,$f(a)$が既知の時,その近傍の点$x=a+h$で
\[
  f(a+h) \approx f(a)+f'(a)(x-a)
\]
が成り立つ.
つまり,\,$f(a)$の値を元に,\,$f(a+h)$を近似したのである.
方程式
\[
  f(x)=x^2-2=0
\]
の正の解 $\sqrt{2}$ を求めたいとする.いまの近似値を $x=a_n$ とし,
点 $(a_n,f(a_n))$ における接線を引く.その接線の傾きは
$f'(a_n)=2a_n$ だから,接線の方程式は
\[
  y=f(a_n)+2a_n(x-a_n)
\]
である.\,$y=0$ と置いて $x$ 軸との交点を求めると,
\begin{align*}
  a_{n+1}
    &=a_n-\frac{f(a_n)}{f'(a_n)}\\
    &=a_n-\frac{a_n^2-2}{2a_n}\\
    &=\frac{a_n^2+2}{2a_n}.
\end{align*}
これが
\[
  a_{n+1}=\frac{a_n^2+2}{2a_n}
\]
の意味である.

例えば $a_1=2$ とすると,
\[
  a_2=\frac{3}{2},
  \quad
  a_3=\frac{17}{12},
  \quad
  a_4=\frac{577}{408}
\]
となり,これらは $\sqrt{2}$ に急速に近づく.この速さは,誤差を
$e_n=a_n-\sqrt{2}$ と置くと見える.
\begin{align*}
  e_{n+1}
    &=\frac{a_n^2+2}{2a_n}-\sqrt{2}\\
    &=\frac{(a_n-\sqrt{2})^2}{2a_n}\\
    &=\frac{e_n^2}{2a_n}.
\end{align*}
誤差がほぼ二乗されるため,正しい近似に近づくと桁数が一気に増える.

この漸化式の解き方は,一般項を閉じた式にすることではない.
「接線を使って,次の近似値をどう設計するか」を考え,その設計が
どのような速さで誤差を減らすかを調べている.付録では,この考えを
漸化式の収束,固定点,誤差評価,微分方程式の数値解法へ広げる.

\subsection{初項依存型}

最後に,漸化式の形だけでなく,初項の選び方そのものが問題を左右する
例を見よう.
\[
  a_1=c,
  \quad
  a_{n+1}=\frac{a_n^2-1}{2}.
\]
初項を記号 $c$ のままにすると,
\begin{align*}
  a_2&=\frac{c^2-1}{2},\\
  a_3&=\frac{c^4-2c^2-3}{8},\\
  a_4&=\frac{c^8-4c^6-2c^4+12c^2-55}{128}.
\end{align*}
一般に $a_n$ は $c$ の次数 $2^{n-1}$ の多項式になる.初項を一般の
文字のまま扱おうとすると,項を一つ進めるたびに次数が倍になり,
三角関数型や基本対称式型のような固定された表現が見えてこない.

特に $c=2$ と固定すると,
\[
  a_1=2,
  \quad
  a_2=\frac{3}{2},
  \quad
  a_3=\frac{5}{8},
  \quad
  a_4=-\frac{39}{128}
\]
となる.しかし,この列は直前の例のように角度を2倍する形にも,
指数の差を保つ形にも,接線近似の形にもならない.

ここでいう「$c=2$ 以外では解けない」とは,他の初項の数列が存在しない
という意味ではない.また,定数列を探すだけなら,
\[
  c=1+\sqrt{2},
  \quad
  c=1-\sqrt{2}
\]
では $a_{n+1}=a_n=c$ となる.大切なのは,ある初項を選んだときだけ
特別な恒等式や置換が働くことがあり,初項を一般化するとその偶然の
構造が失われる,という事実である.

この例では,少なくとも本節で扱った初等的な道具からは一般項を得る
決まった方法が出てこない.それでも漸化式は数列を定めており,
\[
  a_n=f^{\,\circ(n-1)}(c),
  \quad
  f(x)=\frac{x^2-1}{2}
\]
と書くことはできる.これは解答の終わりではなく,「一般項を求める」
という目標自体を見直す入口である.付録では,閉じた式が得られない
場合にも,軌道,極限,近似,行列,母関数など別の言葉で数列を理解する
方法を考える.
